A comparação entre A0 e A1 mostra que a filtragem FIR não deve ser tratada como
sinônimo de melhora na detecção. A remoção de componentes lentas e de
interferência pode tornar a morfologia mais regular, mas também altera sinais
estatísticos que alguns detectores exploram. O One-Class SVM perdeu PR-AUC após
a filtragem, enquanto Isolation Forest e LOF melhoraram. Essa divergência indica
que a qualidade visual do sinal e a separabilidade para um detector one-class
não são propriedades equivalentes.

A entrada de descritores espectrais em A2 favoreceu principalmente Isolation
Forest. Uma interpretação plausível é que energia por bandas, frequência
dominante e estatísticas de STFT criam partições úteis para árvores aleatórias,
sobretudo quando anomalias se afastam por concentração de energia em faixas
específicas. O mesmo acréscimo reduziu o PR-AUC do LOF em relação a A1, o que
sugere perda de estrutura local útil ou aumento de dispersão em dimensões que
não contribuem igualmente para a vizinhança.

As respostas Gabor em A3 produziram o melhor resultado preliminar apenas para
LOF. O padrão é coerente com a hipótese de que energia tempo-frequência
localizada em P, QRS e T pode reorganizar vizinhanças locais de batimentos
normais e anômalos. A mesma representação prejudicou One-Class SVM e Isolation
Forest em relação aos seus melhores estágios, o que recomenda seleção de
features, redução dimensional ou ajuste de hiperparâmetros antes de qualquer
conclusão mais forte sobre Gabor.

A distância entre PR-AUC e métricas em limiar fixo também merece cautela. LOF
atingiu recalls altos com baixa precision, sinalizando grande número de falsos
positivos sob o limiar definido por contaminação. A análise por curva
precision-recall oferece uma leitura mais estável do ranqueamento de anomalias,
enquanto a matriz de confusão depende diretamente de uma escolha operacional de
limiar.

A rotulagem binária normal versus anomalia simplifica classes clínicas com
origens distintas, e as janelas P, QRS e T usadas no protocolo são regiões
aproximadas em torno do R-peak, não delineações clínicas. Ainda assim, a
separação inter-paciente, a exclusão de paced no v1 e o treino restrito a
batimentos normais tornam a primeira execução adequada para uma pergunta de DSP
experimental. O próximo refinamento deve testar sensibilidade a limiares,
seleção de features Gabor e inspeção visual de falsos positivos e falsos
negativos.
